\chapter{Projekt systemu}
\label{ch:projekt}

Niniejszy rozdział przedstawia szczegółowy projekt architektury systemu do przetwarzania i~archiwizacji strumieniowych
danych finansowych.
System został zaprojektowany jako rozproszona architektura mikrousługowa, realizująca paradygmat przetwarzania
sterowanego zdarzeniami (\textit{Event-Driven Architecture}, EDA).
Poszczególne komponenty komunikują się asynchronicznie poprzez centralnego brokera wiadomości, co zapewnia luźne
powiązanie między usługami oraz odporność na chwilowe niedostępności poszczególnych modułów.
Architektura ta umożliwia niezależne skalowanie, wdrażanie i~aktualizowanie każdego z~komponentów, co jest kluczowe
w~kontekście ciągłego rozwoju systemu~\cite{edge_computing_satya}.

Poniżej przedstawiono szczegółowy opis każdego z~komponentów systemu wraz z~opisem przepływu danych między nimi.


\section{Serwer strumieniowy (Data Fetcher)}
\label{sec:data_fetcher}

Serwer strumieniowy stanowi punkt wejścia systemu i~odpowiada za nawiązywanie oraz utrzymywanie połączeń WebSocket
z~interfejsem API giełdy Binance.
Komponent ten został zaimplementowany w~języku Go, który charakteryzuje się niskim narzutem pamięciowym oraz natywnym
wsparciem dla współbieżności poprzez mechanizm gorutyn i~kanałów~\cite{go_at_google}.

Serwer nasłuchuje strumieni danych dla skonfigurowanej listy par handlowych (np.\ BTCUSDT, ETHUSDT) i~odbiera
komunikaty w~czasie rzeczywistym.
Każdy odebrany komunikat jest walidowany pod kątem poprawności struktury JSON, uzupełniany o~znacznik czasowy odbioru
po stronie systemu (dla celów analizy opóźnień), a~następnie publikowany do centralnego brokera wiadomości.

W~przypadku utraty połączenia z~giełdą, serwer realizuje algorytm ponownego łączenia.
Zdarzenia związane z~cyklem życia połączeń (nawiązanie, utrata, błędy) są rejestrowane w~logach z~odpowiednimi
poziomami ważności.

Przepływ danych: \texttt{Binance API} $\rightarrow$ \texttt{Data Fetcher} $\rightarrow$ \texttt{NATS JetStream}


\section{Broker wiadomości (NATS JetStream)}
\label{sec:nats_jetstream}

NATS JetStream pełni rolę centralnej szyny danych (\textit{data bus}) łączącej wszystkie komponenty systemu.
Jest to rozproszony system kolejkowania wiadomości charakteryzujący się wysoką przepustowością, niskimi opóźnieniami oraz
wbudowanym mechanizmem trwałości (\textit{persistence}).

Broker realizuje wzorzec publikuj-subskrybuj (\textit{publish-subscribe}), w~którym producenci (np.\ serwer
strumieniowy) publikują komunikaty do nazwanych strumieni (\textit{streams}), a~konsumenci (np.\ moduł MongoDB)
subskrybują te strumienie i~przetwarzają komunikaty asynchronicznie.
JetStream zapewnia buforowanie komunikatów w~przypadku chwilowej niedostępności konsumentów, gwarantując dostarczenie
po ich wznowieniu działania~\cite{nats_jetstream_docs_persistence_at_least_once}.

Kluczowe cechy konfiguracji JetStream w~systemie:
\begin{itemize}
    \item \textbf{Trwałość} -- komunikaty są zapisywane na dysku, co zapewnia ich przetrwanie w~przypadku restartu
    brokera.
    \item \textbf{Retencja} -- konfigurowalny czas przechowywania komunikatów w~strumieniu (domyślnie 24~godziny).
    \item \textbf{Semantyka dostarczania} -- system wspiera semantykę \textit{at-least-once}, gwarantującą dostarczenie
    każdego komunikatu co najmniej raz~\cite{nats_jetstream_docs_persistence_at_least_once,at_least_once}.
    Konsumenci muszą być idempotentni, aby obsłużyć potencjalne duplikaty.
    \item \textbf{Grupy konsumenckie} (\textit{consumer groups})-- możliwość definiowania grup konsumentów, w~ramach których każdy komunikat
    przetwarzany jest przez dokładnie jednego członka grupy (równoważenie obciążenia).
\end{itemize}


\section{Warstwa gorąca (MongoDB)}
\label{sec:mongodb}

MongoDB pełni rolę operacyjnej bazy danych przechowującej dane bieżące w~warstwie gorącej.
Wybór dokumentowej bazy danych podyktowany jest elastycznością schematu (komunikaty z~giełdy mogą mieć zmienną
strukturę) oraz natywnym wsparciem dla formatu JSON~\cite{mongo_vs_relational}.

System zapewnia automatyczne usuwanie dokumentów po upływie okresu retencji, co odciąża warstwę operacyjną bez
konieczności manualnej interwencji.
Przed usunięciem, dane są przenoszone do warstwy zimnej przez dedykowany proces archiwalny.

Przepływ danych: \texttt{NATS JetStream} $\rightarrow$ \texttt{MongoDB} $\rightarrow$ \texttt{MinIO}


\section{Moduły ETL (Extract, Transform, Load)}
\label{sec:etl_modules}

Moduły ETL odpowiadają za transformację surowych komunikatów z~warstwy gorącej do postaci znormalizowanej oraz ich
zapis do warstwy zimnej (MinIO).
Architektura ETL została podzielona na dwa równoległe potoki:

\begin{enumerate}
    \item \textbf{Potok surowych danych (Raw Writer)} -- zapisuje oryginalne komunikaty bez modyfikacji,
    tworząc archiwum źródłowe służące jako podstawę wiarygodności (\textit{source of truth}).
    Dane surowe umożliwiają debugowanie oraz ponowną transformację w~przypadku zmian w~logice przetwarzania.
    \item \textbf{Potok danych oczyszczonych (Curated Writer)} -- realizuje pełną transformację obejmującą
    walidację strukturalną, normalizację typów danych, wzbogacanie o~metadane oraz obliczanie statystyk agregatowych.
\end{enumerate}

Proces transformacji w~potoku oczyszczonym obejmuje:
\begin{itemize}
    \item Parsowanie komunikatu JSON i~ekstrakcję pól kluczowych (cena, wolumen, znacznik czasowy giełdy, symbol).
    \item Walidację typów i~zakresów wartości (np.\ cena $> 0$, wolumen $\geq 0$).
    \item Obliczenie identyfikatora partycji czasowej (rok/miesiąc/dzień/godzina).
\end{itemize}


\section{Warstwa zimna (MinIO Object Storage)}
\label{sec:minio}

MinIO stanowi warstwę zimną systemu, przeznaczoną do długoterminowego przechowywania danych historycznych.
Jest to magazyn obiektowy kompatybilny z~protokołem Amazon S3, co umożliwia wykorzystanie szerokiego ekosystemu
narzędzi analitycznych~\cite{minio_s3}.

Dane w~MinIO grupowane są hierarchicznie:

Partycjonowanie w~stylu Hive (\texttt{key=value}) umożliwia efektywne przeszukiwanie zbiorów danych przez narzędzia
takie jak Apache Spark, czy Pandas, które mogą wykorzystać strukturę katalogów do eliminacji niepotrzebnych
partycji (\textit{partition pruning})~\cite{hive_partitioning}.


\section{Interfejs GraphQL}
\label{sec:graphql}

Interfejs GraphQL stanowi główny punkt dostępu do danych bieżących przechowywanych w~warstwie gorącej.
W~odróżnieniu od tradycyjnych interfejsów REST, GraphQL umożliwia klientom precyzyjne definiowanie struktury
odpowiedzi, co eliminuje problem nadmiarowego lub niedoborowego pobierania danych (\textit{over-fetching /
under-fetching})~\cite{graphql_rest_comparative}.

Implementacja oparta jest na bibliotece (\textit{gqlgen}), która generuje kod serwera Go na podstawie definicji schematu
GraphQL, zapewniając typowanie statyczne i~wysoką wydajność.


\section{Serwis alertów (Telegram Bot)}
\label{sec:telegram_bot}

Bot platformy Telegram realizuje funkcję powiadamiania użytkowników o~zdarzeniach rynkowych w~czasie rzeczywistym.
Komponent ten subskrybuje strumień NATS JetStream i~filtruje komunikaty pod kątem zdefiniowanych przez użytkowników
reguł alertów.

Przepływ danych: \texttt{NATS JetStream} $\rightarrow$ \texttt{Alert Service} $\rightarrow$ \texttt{Telegram API}
$\rightarrow$ \texttt{Użytkownik}


\section[Analityka strumieniowa i historyczna]{Analityka strumieniowa i~historyczna}
\label{sec:analityka}

System zapewnia dwa tryby dostępu analitycznego:

\begin{enumerate}
    \item \textbf{Analityka w czasie zbliżonym do rzeczywistego (NRT Analytics)} -- realizowana na danych
    przepływających przez system.
    Obejmuje obliczanie agregatów w~oknach czasowych, detekcję trendów oraz identyfikację anomalii.
    \item \textbf{Analityka historyczna (Historical Analytics)} -- realizowana na danych archiwalnych
    przechowywanych w~MinIO\@.
    Dostęp zapewniany jest poprzez środowisko Jupyter Notebook z~bibliotekami Python (Pandas, PyArrow, scikit-learn),
    umożliwiając eksploracyjną analizę danych, budowanie modeli predykcyjnych oraz generowanie raportów.
\end{enumerate}

